\summarychapter{20}{The Nectar of Learning}{Scholasticism}

\begin{inbrief}
As the crusades sought to extend God's rule over Muslim lands, the new universities sought to
understand God's truth wherever it could be found. The chapter describes the schools and their great
teachers, the method called Scholasticism, and the two services the universities rendered the papal
monarchy: a body of canon law and, in Thomas Aquinas, a theology that joined reason and revelation in
one system.
\question{How was the life of the mind brought under God's rule, and reason put to work for faith?}
\end{inbrief}

The Gothic cathedral moves in two directions. Pillars, arches and steeples point upward like man
aspiring to heaven; through the great windows of colored glass, made possible by the outside
buttresses, light descends to the lowly, telling how God saves, from the first day of creation to the day of judgment. Light is the Bible's image of God and of Christ, ``the light of the world.'' In stone the building says what the schoolmen said in words about reason and revelation: man reaches up while God comes down.

\ssection{The Rise of Universities}
Fittingly, the universities grew out of cathedral schools. The period is called Scholasticism for its
distinctive method and its distinctive theology. The Schoolmen had two aims: to reconcile Christian
doctrine with human reason, and to arrange the church's teaching in an orderly system. A free search
for truth was not intended, since the chief doctrines were fixed; discussion was to show that they are
reasonable and to draw out their implications. Curiosity drove it on. Olbert of Li\`ege, in the
eleventh century, flew to every famous master and returned like a bee ``gorged with the nectar of
learning'' to live ``religiously in a studious manner.''

The aim of medieval education was salvation, and at first it was mostly for clergy. Benedict required
his monks to read; Charlemagne ordered every monastery to keep a school and set up a palace school.
The best opportunity for laymen came from the cathedral schools in the towns, founded to train parish
priests and gradually opened to all. They taught the seven liberal arts---grammar, rhetoric, logic,
arithmetic, geometry, music and astronomy---from a few texts: Cassiodorus's \emph{Handbook of
Sacred and Secular Learning}, the \emph{Consolation of Philosophy} that Boethius wrote in prison
under King Theodoric, Augustine, Gregory the Great and a few other fathers. These were the
\emph{auctores}, authorities not lightly contradicted.

\ssection{The Magnetism of an Able Teacher}
Universities were born from the drawing power of great teachers. The first, as Anne Fremantle
describes them, was Gerbert, who led the school attached to the cathedral of Rheims shortly before the year 1000, a
serf's son who died as Pope Sylvester II (999--1003). Sent as a young monk to study mathematics in
Spain, he met the inquiring spirit of Muslim scholarship, the first of several Muslim contributions to
the Christian awakening. At Rheims he declared that quoting the authorities was not enough, had his
pupils read the Roman classics in the original, and collected a library at a time when a manuscript
took a year to copy.

The liveliest figure was Peter Abelard (1079--1142), eldest son of a minor noble of Brittany, who gave
up his inheritance for learning, wandered among the masters, challenged them in class and became a
celebrated lecturer in Paris. His \emph{Sic et Non} (``Yes and No'') set out 158 questions and
answered each with conflicting quotations from Scripture, the fathers and the pagan classics. ``By
doubting we come to inquiry, and by inquiry we arrive at the truth.'' A council at Soissons condemned
his book on the Trinity in 1121 and he was confined to a monastery. A year later he was allowed to
live in the wilderness southeast of Paris, where students flocked to him, built him a shelter and
tilled his land. He fell foul of Bernard of Clairvaux, for whom ``the faith of the righteous
believes; it does not dispute.'' At Bernard's urging a council at Sens condemned him for heresy in
1140, and he spent his last two years at Cluny.

The seed he had sown could not be stifled. Within a century universities flourished at Paris,
Orleans and Montpellier, at Oxford and Cambridge, at Bologna and Padua. Their flowering was marked by
the grouping of scholars into guilds, each called a \emph{universitas}, the name for any corporate
body. In Italy, where the students were mature men studying law and medicine, the students' guilds
hired and paid the teachers and fined lecturers who skipped a chapter. In France and England, where
students were younger, the masters ruled, forbidding swearing and gambling and prescribing table
manners. There were no buildings: lectures were given in sheds, cloisters and squares, later in rented
rooms with straw on the floor, and a university at odds with the townspeople could simply move.
Besides the lecture the method was the disputation, in which masters debated texts by Abelard's
questions and answers. Scholasticism came to mean the painstaking reaching of logical conclusions by
questioning, examining and arranging. Wars of logic ran for years amid stamping and whistling, but
students were learning to think, and even Christian doctrine was being examined.

\ssection{Support for the Papal Monarchy}
All this went on under the eye of the papacy, which saw to it that the outcome supported the papal
monarchy. The first support was canon law, the church's law, administered in every church court from
the bishop's to the pope's. Its center was Bologna. About 1140 Gratian, a monk there, published his
\emph{Harmony of Discordant Canons}, arranging earlier collections by subject. It became the
standard manual and the base of the Body of Canon Law, which stayed in force until its revision in
1918. Canon law intervened in wars and protected widows, fixed fasts and feasts, required confession
and communion at least once a year and defined what brought excommunication. Through baptism and
marriage it reached the most intimate relations of men and women, and it bound plowman and prince
alike. Shelley sees in this zeal to control every part of life a Christian legalism that shrank human
freedom to a minimum and forgot that freedom prepares the way for true faith. But canon law gave the
papacy a rational legal basis that no medieval state yet had, and with it an international standing
above any kingdom. Innocent III preached at his consecration on Jeremiah 1:10, ``I have this day set
thee over the nations,'' and acted on it.

The second support was a rational theology of Christian society. In the thirteenth century lost
works of Aristotle, preserved among eastern Christians and especially Muslims, returned to the West.
Theologians since Origen and Augustine had assumed Plato's framework: an eternal realm of perfect,
immaterial ideas and below it a material world of imperfect, perishable copies. Aristotle held that
the forms are present in concrete things and are found by studying the created world, so that thought
and matter fit together. Many knew him only through the Muslim commentator Averroes, who stressed the
independence of reason from revelation and followed Aristotle to conclusions that contradict
Christian teaching, such as the eternity of matter. Thomas Aquinas's life work was to insist that
sound thinking never leads away from God's truth. Revelation surpasses what reason can reach but
never contradicts it: ``All truth is God's truth.''

\ssection{The Peak of Rational Reach}
Thomas Aquinas (1224--1274), a Dominican of noble birth from Aquino, brilliant, tireless and gentle,
was sent from Italy to Paris to deal with these controversies. His \emph{Summa Theologica} treats the
whole universe from the standpoint of God. Philosophy and theology, reason and revelation, are
distinct but both come from God. Reason starts from the visible creation and reaches ``the vestibule
of faith.'' It can prove God's existence, for the chain of causes leads back to a First Cause.
Revelation, which is more certain, gives the full knowledge of God, such as the Trinity, and of man's
origin and destiny.

Man is a sinner in need of grace. Christ's sacrifice won reconciliation; the church applies it.
Christians need a constant infusion of ``cooperating grace'' that stirs the virtues, above all love,
and enables works that earn merit. This grace comes only through the sacraments entrusted to the
visible Roman church under the pope, submission to whom Aquinas held necessary for salvation. Following
Peter Lombard, author of the standard textbook, he counted seven sacraments (baptism, confirmation, penance, the Lord's Supper, marriage, ordination and extreme unction). The greatest is the Supper, a
true sacrifice continuing that of the cross, in which the substance of bread and wine is changed into
the body and blood of Christ while the appearance remains---transubstantiation, of which Thomas
gave the classic statement. Penance, the sacrament of healing, consists of contrition, confession to
the priest and satisfaction. With some cautions he accepted indulgences: the merits of Christ and the
saints form a ``treasury of merit'' on which priests may draw for those with too little of their own.
At death the wicked go to hell and the faithful who have used the means of grace to heaven; the great
mass must be purified in purgatory, where the prayers of the church can relieve them.

None of this was new, but Thomas set it in a grand, almost cosmic framework that, like the cathedral,
harmonized human aspiration and divine light and gave the papal monarchy a supernatural view of this
world and the next. Here, Shelley says, is the peak of ecclesiastical aspiration and arrogance: pope
and priests reaching beyond the grave. Like the crusades, Scholasticism may have claimed too much.
Against Innocent's boast to stand ``below God and above man,'' a growing number of unlettered laymen
remembered the apostle's words, ``There is one God and one Mediator.''

\begin{keyfacts}
\keypeople{Gerbert (Sylvester II); Peter Abelard; Bernard of Clairvaux; Cassiodorus; Boethius;
Gratian; Aristotle; Averroes; Thomas Aquinas; Peter Lombard; Innocent III.}
\keydates{%
  \kd{999--1003}{Gerbert of Rheims reigns as Pope Sylvester II}
  \kd{1079--1142}{Peter Abelard}
  \kd{1121}{Council of Soissons condemns Abelard's book on the Trinity}
  \kd{1140}{Council of Sens condemns Abelard}
  \kd{c. 1140}{Gratian's \emph{Harmony of Discordant Canons}}
  \kd{1224--1274}{Thomas Aquinas, author of the \emph{Summa Theologica}}}
\keyterms{Scholasticism; Schoolmen; seven liberal arts; \emph{auctores}; \emph{universitas};
disputation; canon law; First Cause; cooperating grace; seven sacraments; transubstantiation;
treasury of merit; purgatory.}
\end{keyfacts}
